% sections/results.tex -- Section 7 (paper_v2_vn, translated in L6 from paper_v2's
% sections/results.tex; terms as in glossary_vn.tex). The details that left it in paper_v2 (the
% paths' ranks, the delta check, the user groups, more paths, the path weights and preferences, HAN
% against metapath2vec) are in the supplementary document (supplement_vn.tex, N31). A few of their
% numbers are given here in the text. Every number is \val{} from numbers.tex. Tables: T4
% (tab:main), T7 (tab:cost); figures: F3 (fig:sets),
% F4 (fig:weights).
\section{Kết quả}
\label{sec:results}

Mọi con số đều là kết quả trên tập thử nghiệm trừ khi có ghi chú khác.

\subsection{Tập tìm được so với tập thủ công}
\label{sec:res-main}

Theo độ đo đã cố định từ trước, \code{a1} thua tập thủ công cùng kích thước trên cả hai mô hình
(Bảng~\ref{tab:main}). Trong HAN nó kém M5 \val{han-4-hit-diffabs} (\val{han-4-hit-relabs}\%,
\val{han-4-hit-pt}) và kém M6 \val{han-6-hit-diffabs} (\val{han-6-hit-relabs}\%,
\val{han-6-hit-pt}), bằng \val{han-4-hit-sdx} và \val{han-6-hit-sdx} lần độ lệch chuẩn giữa các
seed (lấy bên lớn hơn). Trong metapath2vec, chênh lệch rộng hơn: \val{mtv-4-hit-diffabs}
(\val{mtv-4-hit-relabs}\%) và \val{mtv-6-hit-diffabs} (\val{mtv-6-hit-relabs}\%). Khi xếp hạng
toàn bộ, các tập của HAN ngang bằng nhau (\val{han-4-fhit-pt} và \val{han-6-fhit-pt}). Trong
metapath2vec, \code{a1} ngang bằng với $M = 4$ meta-path (\val{mtv-4-fhit-pt}) và thua với $M = 6$ (\val{mtv-6-fhit-pt}). Trong
HAN sự thua kém này cũng đúng với \code{a2} và \code{c} (Hình~\ref{fig:sets}) và ở mọi nhóm người
dùng (chênh lệch từ \val{han-6-hit-d3-diff} đến \val{han-4-hit-d1-diff}).

% t4_main.tex -- T4: the main comparison (paper_v2_vn, translated in L3 from paper_v2's
% tables/t4_main.tex; results.csv via numbers.tex; Welch tests computed by make_numbers.py).
\begin{widetable}[htbp]
  \centering
  \small
  \setlength{\tabcolsep}{4pt}
  \caption{Tập tìm được so với tập thủ công: kết quả trên tập thử nghiệm, trung bình $\pm$ độ lệch
    chuẩn qua 3 seed huấn luyện. $\Delta$: tập tìm được trừ tập thủ công, kèm $p$ của kiểm định
    Welch hai phía. Lấy mẫu: mỗi cạnh thử nghiệm được xếp hạng so với \val{split-negatives} mẫu âm,
    với Hit@10 là độ đo đã cố định từ trước. Toàn bộ: so với mọi sản phẩm mà người dùng không có
    cạnh nối tới, với Recall@20 lấy trung bình theo người dùng.}
  \label{tab:main}
  \begin{tabular*}{\linewidth}{@{\extracolsep{\fill}}llllll@{}}
    \toprule
    & & \multicolumn{2}{c}{Lấy mẫu} & \multicolumn{2}{c}{Xếp hạng toàn bộ} \\
    \cmidrule(lr){3-4}\cmidrule(l){5-6}
    Mô hình & Tập & Hit@10 & NDCG@10 & Hit@10 & Recall@20 \\
    \midrule
    HAN & \code{a1}, 4 meta-path & \val{han-f4a1-hit-mean} $\pm$ \val{han-f4a1-hit-sd}
      & \val{han-f4a1-ndcg-mean} & \val{han-f4a1-fhit-mean} $\pm$ \val{han-f4a1-fhit-sd}
      & \val{han-f4a1-frec-mean} \\
    & M5 & \val{han-m5-hit-mean} $\pm$ \val{han-m5-hit-sd} & \val{han-m5-ndcg-mean}
      & \val{han-m5-fhit-mean} $\pm$ \val{han-m5-fhit-sd} & \val{han-m5-frec-mean} \\
    & $\Delta$ ($p$) & \val{han-4-hit-diff} (\val{han-4-hit-p})
      & \val{han-4-ndcg-diff} (\val{han-4-ndcg-p}) & \val{han-4-fhit-diff} (\val{han-4-fhit-p})
      & \val{han-4-frec-diff} (\val{han-4-frec-p}) \\
    \addlinespace
    & \code{a1}, 6 meta-path & \val{han-f6a1-hit-mean} $\pm$ \val{han-f6a1-hit-sd}
      & \val{han-f6a1-ndcg-mean} & \val{han-f6a1-fhit-mean} $\pm$ \val{han-f6a1-fhit-sd}
      & \val{han-f6a1-frec-mean} \\
    & M6 & \val{han-m6-hit-mean} $\pm$ \val{han-m6-hit-sd} & \val{han-m6-ndcg-mean}
      & \val{han-m6-fhit-mean} $\pm$ \val{han-m6-fhit-sd} & \val{han-m6-frec-mean} \\
    & $\Delta$ ($p$) & \val{han-6-hit-diff} (\val{han-6-hit-p})
      & \val{han-6-ndcg-diff} (\val{han-6-ndcg-p}) & \val{han-6-fhit-diff} (\val{han-6-fhit-p})
      & \val{han-6-frec-diff} (\val{han-6-frec-p}) \\
    \midrule
    metapath2vec & \code{a1}, 4 meta-path & \val{mtv-f4a1-hit-mean} $\pm$ \val{mtv-f4a1-hit-sd}
      & \val{mtv-f4a1-ndcg-mean} & \val{mtv-f4a1-fhit-mean} $\pm$ \val{mtv-f4a1-fhit-sd}
      & \val{mtv-f4a1-frec-mean} \\
    & M5 & \val{mtv-m5-hit-mean} $\pm$ \val{mtv-m5-hit-sd} & \val{mtv-m5-ndcg-mean}
      & \val{mtv-m5-fhit-mean} $\pm$ \val{mtv-m5-fhit-sd} & \val{mtv-m5-frec-mean} \\
    & $\Delta$ ($p$) & \val{mtv-4-hit-diff} (\val{mtv-4-hit-p})
      & \val{mtv-4-ndcg-diff} (\val{mtv-4-ndcg-p}) & \val{mtv-4-fhit-diff} (\val{mtv-4-fhit-p})
      & \val{mtv-4-frec-diff} (\val{mtv-4-frec-p}) \\
    \addlinespace
    & \code{a1}, 6 meta-path & \val{mtv-f6a1-hit-mean} $\pm$ \val{mtv-f6a1-hit-sd}
      & \val{mtv-f6a1-ndcg-mean} & \val{mtv-f6a1-fhit-mean} $\pm$ \val{mtv-f6a1-fhit-sd}
      & \val{mtv-f6a1-frec-mean} \\
    & M6 & \val{mtv-m6-hit-mean} $\pm$ \val{mtv-m6-hit-sd} & \val{mtv-m6-ndcg-mean}
      & \val{mtv-m6-fhit-mean} $\pm$ \val{mtv-m6-fhit-sd} & \val{mtv-m6-frec-mean} \\
    & $\Delta$ ($p$) & \val{mtv-6-hit-diff} (\val{mtv-6-hit-p})
      & \val{mtv-6-ndcg-diff} (\val{mtv-6-ndcg-p}) & \val{mtv-6-fhit-diff} (\val{mtv-6-fhit-p})
      & \val{mtv-6-frec-diff} (\val{mtv-6-frec-p}) \\
    \midrule
    Phổ biến & -- & \val{pop-hit} & \val{pop-ndcg} & \val{pop-fhit} & \val{pop-frec} \\
    \bottomrule
  \end{tabular*}
\end{widetable}


\subsection{Đối chứng ngẫu nhiên}
\label{sec:res-random}

Để biết khoảng cách giữa tìm kiếm và thiết kế thủ công lớn đến đâu so với ảnh hưởng của việc chọn
meta-path, Hình~\ref{fig:sets} đặt các tập tìm được và các tập thủ công giữa cả 24 tập, cùng với các tập
chọn ngẫu nhiên làm đối chứng. Trong tám tập của mỗi kích thước được lấy
từ cùng \val{search-cands} ứng viên, gồm ba tập tìm được và năm tập ngẫu nhiên, HAN xếp các tập
tìm được ở vị trí \val{han-4-pos-hit} với $M = 4$ meta-path và \val{han-6-pos-hit} với $M = 6$, giống như
BPR xác thực. Khi xếp hạng toàn bộ, các tập này đứng ở vị trí \val{han-4-pos-fhit} với $M = 4$ và
\val{han-6-pos-fhit} với $M = 6$. Các khoảng cách này nằm trong nhiễu giữa các seed: tập tìm được tốt
nhất đạt \val{han-f4c-hit-s1} với $M = 4$ và \val{han-f6a1-hit-s1} với $M = 6$, so với
\val{han-r4-hit-min} và \val{han-r6-hit-min} của những tập ngẫu nhiên yếu nhất. Vì vậy, các vị trí này cho
thấy tìm kiếm không vượt được chính các ứng viên của nó, chứ không cho thấy mọi tập ngẫu nhiên đều vượt nó.
Các tập thủ công đạt kết quả xấp xỉ một tập ngẫu nhiên trung bình: M5 đạt \val{han-m5-hit-mean} so với
trung bình \val{han-r4-hit-mean} của các tập ngẫu nhiên $M = 4$ meta-path (khoảng
\val{han-r4-hit-min}--\val{han-r4-hit-max}) và M6 đạt \val{han-m6-hit-mean} so với
\val{han-r6-hit-mean} (\val{han-r6-hit-min}--\val{han-r6-hit-max}).

% f3_sets.tex -- F3: HAN test Hit@10 of all 24 sets by family (paper_v2_vn, translated in L3 from
% paper_v2's figures/f3_sets.tex). Data: f3_*.dat, written by make_numbers.py from results.csv
% (seed 1 of every set).
\begin{widefigure}[htbp]
  \centering
  \begin{tikzpicture}
    \begin{axis}[
        width=0.9\linewidth, height=5.4cm,
        xmin=0.45, xmax=3.6, ymin=0.612, ymax=0.712,
        xtick={1,2,3}, xticklabels={2--3 meta-path, 4 meta-path, 6 meta-path},
        ylabel={Hit@10 thử nghiệm của HAN (lấy mẫu)},
        yticklabel style={/pgf/number format/fixed, /pgf/number format/precision=2},
        ytick={0.62,0.64,0.66,0.68,0.70},
        tick label style={font=\scriptsize}, label style={font=\footnotesize},
        ymajorgrids, grid style={black!8}, axis line style={black!60}, tick style={black!60},
        legend style={font=\scriptsize, at={(0.02,0.98)}, anchor=north west, draw=none,
          fill=white, fill opacity=0.85, text opacity=1},
        legend cell align=left]
      \draw[black!20] (axis cs:1.5,0.612) -- (axis cs:1.5,0.712);
      \draw[black!20] (axis cs:2.5,0.612) -- (axis cs:2.5,0.712);
      \addplot[only marks, mark=*, mark size=2.6pt, mark options={fill=cSearch, draw=white,
          line width=0.4pt}] table[x=x, y=hit] {figures/f3_searched.dat};
      \addlegendentry{tìm được (LMSPS)}
      \addplot[only marks, mark=square*, mark size=2.4pt, mark options={fill=cManual, draw=white,
          line width=0.4pt}] table[x=x, y=hit] {figures/f3_manual.dat};
      \addlegendentry{thủ công}
      \addplot[only marks, mark=triangle*, mark size=3pt, mark options={fill=cRandom, draw=white,
          line width=0.4pt}] table[x=x, y=hit] {figures/f3_random.dat};
      \addlegendentry{chọn ngẫu nhiên}
      \addplot[only marks, mark=none, point meta=explicit symbolic,
          visualization depends on={value \thisrow{anchor}\as\labelanchor},
          nodes near coords, nodes near coords style={font=\scriptsize, text=black!80,
            anchor=\labelanchor, inner sep=2.5pt}]
        table[x=x, y=y, meta=label] {figures/f3_labels.dat};
    \end{axis}
  \end{tikzpicture}
  \Description{Biểu đồ điểm Hit@10 thử nghiệm của HAN cho cả 24 tập meta-path, nhóm theo kích thước tập, mỗi loại tập một hình dạng và màu điểm riêng. Các tập tìm được nằm ở phía thấp của các tập lấy từ các ứng viên. Các tập thủ công nằm giữa các tập chọn ngẫu nhiên.}
  \caption{Hit@10 thử nghiệm của HAN cho cả 24 tập meta-path, mỗi tập một seed (độ lệch chuẩn qua
    ba seed của các tập chính nhiều nhất là \val{han-m6-hit-sd}). Ở seed này, các tập tìm được đứng
    cuối trong số các tập lấy từ cùng \val{search-cands} ứng viên ở cả hai kích thước, dù với chênh
    lệch nằm trong nhiễu giữa các seed. M5 và M6 nằm gần giữa các tập chọn ngẫu nhiên. Để so sánh:
    mô hình phổ biến \val{pop-hit}, GAT \val{gat-hit-mean}.}
  \label{fig:sets}
\end{widefigure}


Do đó, việc chọn meta-path ít ảnh hưởng đến HAN: cả \val{han-4-nsets} tập $M = 4$ meta-path chỉ chênh
nhau trong phạm vi \val{han-4-span} Hit@10, còn cả \val{han-6-nsets} tập $M = 6$ meta-path chênh nhau trong phạm vi
\val{han-6-span}. Chọn theo xác thực là hiệu quả: trên 24 tập, BPR xác thực và Hit@10 thử nghiệm có
hệ số tương quan Spearman \val{han-valtest-rho} (\val{han-valfull-rho} với Hit@10 xếp hạng toàn
bộ). Tập chọn theo xác thực vốn là một tập ngẫu nhiên đạt \val{han-r6d2-hit-s1}, mức mà không tập
nào khác vượt được theo độ đo đã cố định từ trước. Như vậy, lựa chọn rẻ nhất là chọn ngẫu nhiên vài tập
rồi giữ tập có điểm xác thực tốt nhất. Trong metapath2vec, so sánh với ngẫu nhiên cho kết quả trái
chiều: các tập tìm được đứng \val{mtv-4-pos-hit} trong tám tập với $M = 4$ meta-path, trên tập ngẫu nhiên
trung bình, nhưng \val{mtv-6-pos-hit} với $M = 6$ meta-path, dưới mức đó. Hai mô hình cũng ưa thích các
tập khác nhau: trên 23 tập có meta-path qua người dùng, giá trị Hit@10 của chúng có hệ số tương
quan Spearman chỉ \val{han-mtv-rho}.

\subsection{Chi phí so với hiệu quả}
\label{sec:res-cost}

Thiết kế thủ công không dùng GPU và thời gian của nó cũng không được ghi lại (Bảng~\ref{tab:cost}).
Các công trình trước cũng chỉ mô tả chi phí này bằng lời \citep[][Phụ lục~D]{han2020genetic,li2024longrange}.
LMSPS còn bỏ thời gian thiết kế thủ công ra khỏi phép so sánh hiệu quả của mình, cùng với thời gian
tìm kiếm của chính nó \citep[Phụ lục~E.1]{li2024longrange}. Tuy vậy, chi phí của tìm kiếm cho chúng
tôi một mốc để so sánh: toàn bộ \val{search-runs-total} lần chạy tìm kiếm mất khoảng
\val{search-total-h}~giờ~GPU, chưa kể thời gian điều chỉnh mã nguồn. Thiết kế thủ công rẻ hơn tìm
kiếm chừng nào việc thiết kế mất ít hơn khoảng thời gian này. Chúng tôi cho rằng điều đó hợp lý trên
một lược đồ chỉ có ba kiểu nút, nơi các tập thủ công được thiết kế theo bốn nguyên tắc trong một
buổi, cùng các phép kiểm tra độ dày đặc chỉ mất vài phút CPU. Trong khi đó, tìm kiếm thu được những
tập kém các tập thủ công \val{han-4-hit-relabs}\% và \val{han-6-hit-relabs}\% theo độ đo đã cố định
từ trước và chỉ ngang bằng chúng khi xếp hạng toàn bộ. Như vậy, trên lược đồ này, khoản chi phí tìm
kiếm chưa đem lại cho HAN kết quả vượt trội hơn thiết kế thủ công. Đối chứng ngẫu nhiên vốn không
tốn chi phí, cũng đạt kết quả ít nhất ngang bằng cả hai cách. Khi xếp hạng toàn bộ, trung bình kết
quả của các tập ngẫu nhiên ít nhất cũng cao bằng, dù phân tán rộng hơn
(\val{han-r4-fhit-min}--\val{han-r4-fhit-max} và \val{han-r6-fhit-min}--\val{han-r6-fhit-max}).

% t7_cost.tex -- T7: cost and performance of the three ways to choose HAN's paths (paper_v2_vn,
% translated in L3 from paper_v2's tables/t7_cost.tex). Search cost from the 60 search logs (150
% epochs at the logged epoch time plus the feature propagation); results from results.csv; all via
% numbers.tex. The cost of the hand design is stated in words: none of the reference papers
% measures it.
\begin{widetable}[htbp]
  \centering
  \small
  \caption{Chi phí và kết quả của từng cách chọn meta-path cho HAN: Hit@10 trên tập thử
    nghiệm, lấy mẫu và xếp hạng toàn bộ. Thủ công và bằng tìm kiếm: trung bình qua 3 seed. Ngẫu
    nhiên: trung bình của 5 lần chọn, mỗi lần một seed. Chênh lệch tương đối được tính so với tập
    thủ công cùng kích thước. Huấn luyện HAN trên tập đã chọn tốn như nhau cho cả ba cách, trung vị
    \val{han-run-min}~phút~GPU mỗi lần chạy.}
  \label{tab:cost}
  \begin{tabularx}{\linewidth}{@{}lLlll@{}}
    \toprule
    & & \multicolumn{2}{l}{Hit@10 lấy mẫu} & Hit@10 toàn bộ \\
    \cmidrule(lr){3-4}\cmidrule(l){5-5}
    Cách chọn & Chi phí lựa chọn & 4 / 6 meta-path & so với thủ công & 4 / 6 meta-path \\
    \midrule
    Thủ công (M5, M6) & không dùng GPU: một buổi thiết kế có AI hỗ trợ (không ghi lại thời gian)
      cùng các phép kiểm tra độ dày đặc trên một đồ thị nhỏ, mất vài phút CPU
      & \val{han-m5-hit-mean} / \val{han-m6-hit-mean} & --
      & \val{han-m5-fhit-mean} / \val{han-m6-fhit-mean} \\
    Bằng tìm kiếm (\code{a1}) & \val{search-runs} lần chạy cho mỗi kích thước, trung vị
      \val{search-run-min-4} / \val{search-run-min-6}~phút~GPU mỗi lần:
      \val{search-total-h}~giờ~GPU & \val{han-f4a1-hit-mean} / \val{han-f6a1-hit-mean}
      & \val{han-4-hit-rel}\% / \val{han-6-hit-rel}\%
      & \val{han-f4a1-fhit-mean} / \val{han-f6a1-fhit-mean} \\
    Ngẫu nhiên (5 lần chọn) & không tốn & \val{han-r4-hit-mean} / \val{han-r6-hit-mean}
      & \val{han-r4-vs-manual-rel}\% / \val{han-r6-vs-manual-rel}\%
      & \val{han-r4-fhit-mean} / \val{han-r6-fhit-mean} \\
    \bottomrule
  \end{tabularx}
\end{widetable}


\subsection{HAN làm gì với các meta-path}
\label{sec:res-weights}

Hình~\ref{fig:weights} so sánh mức ưa thích của tìm kiếm đối với từng meta-path với trọng số mà
HAN gán cho nó. Hai bên thống nhất về các meta-path qua người dùng. \code{PUPUP} vốn nối các sản
phẩm qua hai người dùng, có mặt trong cả \val{searched-with-pupup} tập tìm được và \code{PUPPUP}
trong \val{searched-with-puppup} tập. Cả hai đều xếp gần đầu trong các ứng viên của tìm kiếm.
HAN không bao giờ loại chúng hay \code{PUP}, đồng thời gán cho cả ba trọng số cao hơn mức chia
đều. Hai bên bất đồng về danh mục. Mỗi tập tìm được chứa \val{searched-empty-min} đến
\val{searched-empty-max} meta-path danh mục gần rỗng mà tìm kiếm xếp hạng cao, nhưng HAN chỉ giữ
một trong số đó cho mỗi tập với trọng số \val{han-empty-kept-min} đến \val{han-empty-kept-max}
(\val{han-empty-kept-f4c} với \code{c} ở $M = 4$ meta-path), nên các tập tìm được $M = 6$ vận hành
như các tập $M = 4$. Các meta-path danh mục dày đặc của các tập thủ công xếp hạng thấp trong
tìm kiếm nhưng nhận được khoảng mức chia đều hoặc hơn trong HAN (trung vị \val{kind-cat-weight}).
Không tập tìm được nào chứa một meta-path danh mục đến được phần lớn sản phẩm.

% f4_weights.tex -- F4: HAN's path weights against the search's preference, by kind of path
% (paper_v2_vn, translated in L3 from paper_v2's figures/f4_weights.tex). Data: f4_*.dat, written
% by make_numbers.py (search records; HAN seed-1 logs).
\begin{figure}[htbp]
  \centering
  \begin{tikzpicture}
    \begin{axis}[
        width=\linewidth, height=5.6cm,
        xmin=0, xmax=33, ymin=-0.1, ymax=2.05, xtick={1,8,16,24,32},
        xlabel={Mức ưa thích của tìm kiếm: hạng $\alpha$ trung bình (1 = tốt nhất trong
          \val{search-cands})},
        ylabel={Trọng số của HAN (1 = mức chia đều)},
        tick label style={font=\scriptsize}, label style={font=\footnotesize},
        ymajorgrids, grid style={black!8}, axis line style={black!60}, tick style={black!60},
        legend style={font=\scriptsize, at={(0.5,1.02)}, anchor=south, draw=none,
          legend columns=2, /tikz/every even column/.append style={column sep=6pt}},
        legend cell align=left]
      \draw[black!45, dashed] (axis cs:0,1) -- (axis cs:33,1);
      \addplot[only marks, mark=*, mark size=2.6pt, mark options={fill=cSearch, draw=white,
          line width=0.4pt}] table[x=rank, y=weight] {figures/f4_users.dat};
      \addlegendentry{qua người dùng (\val{kind-users-paths} meta-path)}
      \addplot[only marks, mark=square*, mark size=2.4pt, mark options={fill=cManual, draw=white,
          line width=0.4pt}] table[x=rank, y=weight] {figures/f4_empty.dat};
      \addlegendentry{danh mục gần rỗng (\val{kind-empty-paths})}
      \addplot[only marks, mark=triangle*, mark size=3pt, mark options={fill=cRandom, draw=white,
          line width=0.4pt}] table[x=rank, y=weight] {figures/f4_cat.dat};
      \addlegendentry{danh mục khác (\val{kind-cat-paths})}
      \addplot[only marks, mark=diamond*, mark size=3pt, mark options={fill=cOther, draw=white,
          line width=0.4pt}] table[x=rank, y=weight] {figures/f4_pp.dat};
      \addlegendentry{chỉ sản phẩm (\val{kind-pp-paths})}
      \addplot[only marks, mark=none, point meta=explicit symbolic,
          visualization depends on={value \thisrow{anchor}\as\labelanchor},
          nodes near coords, nodes near coords style={font=\scriptsize\ttfamily, text=black!80,
            anchor=\labelanchor, inner sep=2.5pt}]
        table[x=x, y=y, meta=label] {figures/f4_labels.dat};
    \end{axis}
  \end{tikzpicture}
  \Description{Biểu đồ phân tán, mỗi điểm là một meta-path: mức ưa thích trung bình của tìm kiếm đối với meta-path so với trọng số trung bình mà HAN gán cho nó. Cả hai đều ưa các meta-path qua người dùng. Các meta-path danh mục gần rỗng được tìm kiếm ưa thích nhưng gần như không có trọng số trong HAN. Các meta-path danh mục dày đặc bị tìm kiếm xếp hạng thấp nhưng được HAN sử dụng.}
  \caption{HAN làm gì với các meta-path: mỗi điểm là một meta-path, đặt theo mức ưa thích trung
    bình của tìm kiếm đối với nó (cả hai kích thước tập) và trọng số trung bình HAN gán cho nó trên
    các tập chứa nó (seed 1). Hai bên thống nhất về các meta-path qua người dùng. Hai bên bất đồng
    về các meta-path danh mục gần rỗng, vốn được tìm kiếm xếp hạng cao còn HAN phần lớn loại bỏ
    (\val{kind-empty-dropped} trong \val{kind-empty-uses} lần dùng có trọng số dưới
    \val{weight-drop}) và về các meta-path danh mục dày đặc, vốn bị tìm kiếm xếp hạng thấp còn HAN
    sử dụng.}
  \label{fig:weights}
\end{figure}


Trong các tập $M = 4$ và $M = 6$ meta-path, sự thiếu vắng này phân biệt các tập tìm được với tất cả các
tập khác. Ở cả hai kích thước, mọi tập có một meta-path danh mục đến được gần như mọi sản phẩm
đều vượt mọi tập không có meta-path như vậy: \val{han-4-dense-min}--\val{han-4-dense-max} so với
\val{han-4-nodense-min}--\val{han-4-nodense-max} với $M = 4$ (cặp gần nhau nhất nằm trong
nhiễu giữa các seed) và \val{han-6-dense-min}--\val{han-6-dense-max} so với \val{han-6-nodense-max}
với $M = 6$. Trên \val{pref-cand-n} tập lấy từ các ứng viên, tìm kiếm càng ưa thích các meta-path của
một tập thì HAN làm càng kém với tập đó (hệ số Spearman \val{pref-cand-rho} giữa Hit@10 thử nghiệm và
hạng $\alpha$ trung bình của các meta-path, trong đó hạng càng thấp nghĩa là càng được ưa thích,
\val{pref-cand-pt}). Quan hệ này dựa trên các tập tìm được: nếu chỉ xét mười tập chọn ngẫu nhiên,
hệ số này chỉ còn \val{pref-random-rho} (\val{pref-random-pt}).

\paragraph{Nơi tập trung sự thua kém.} Một bước đơn lẻ trên cây phân cấp danh mục, \code{PCCP} nối
\val{bg-pccp-reach-products-pct}\% số sản phẩm với một sản phẩm khác, phần lớn là DVD và video vốn
chiếm \val{bg-pccp-reach-test-pct}\% số cạnh thử nghiệm. Trên các sản phẩm mà \code{PCCP} không đến
được, các tập tìm được kém các tập có meta-path danh mục dày đặc \val{bg-4-gapabs-no} Hit@10 với $M = 4$ meta-path
và \val{bg-6-gapabs-no} với $M = 6$ meta-path (seed~1). Trên các sản phẩm nó đến được, các tập tìm được chỉ kém
\val{bg-4-gapabs-reach} và \val{bg-6-gapabs-reach}. Như vậy, sự thua kém chủ yếu nằm ở sách và
nhạc, những sản phẩm mà các meta-path danh mục của các tập tìm được là rỗng.

\paragraph{Vì sao các tập tìm được thua.} Phép kiểm tra trên chỉ ra các meta-path danh mục gần
rỗng. Trong tìm kiếm, một sản phẩm mà meta-path như vậy không đến được nhận một đặc trưng toàn số 0
cho meta-path đó. Số 0 ấy vẫn mang thông tin: các sản phẩm mà một bước trên cây phân cấp đến được,
phần lớn là DVD và video, chiếm \val{pccp-reach-train-pct}\% số cạnh huấn luyện nhưng chỉ
\val{bg-pccp-reach-products-pct}\% số sản phẩm và do đó cũng chừng ấy phần các mẫu âm ngẫu nhiên,
nên các meta-path rỗng làm lộ nhóm sản phẩm và một phần độ phổ biến của sản phẩm. Trong HAN, tín
hiệu này phần lớn biến mất: mỗi sản phẩm giữ đặc trưng riêng của nó qua một cạnh tự nối trong mọi
đồ thị, nên một kênh rỗng chỉ lặp lại các đặc trưng đó. HAN loại bỏ phần lớn các kênh này. Thứ mà
các tập tìm được khi đó còn thiếu, một meta-path danh mục đến được sách và nhạc, chính là chỗ chúng
thua nhiều nhất. Tìm kiếm cũng có thể tận dụng một meta-path dày đặc kém hơn HAN, vì nó lấy trung
bình đặc trưng của meta-path trên mọi điểm cuối, trong khi HAN chỉ chú ý trên \val{han-set-topk} láng
giềng được chọn. Tuy nhiên, phép kiểm tra theo nhóm sản phẩm chỉ xác định được nơi xảy ra sự thua
kém chứ chưa cho thấy một meta-path danh mục dày đặc sẽ loại bỏ được nó.

\subsection{Các mô hình khác}
\label{sec:res-models}

GAT vượt cả hai mô hình dùng meta-path trên mọi độ đo và ở mọi nhóm người dùng. Nó đạt Hit@10
\val{gat-hit-mean}, so với
\val{han-r6d2-hit-s1} của tập chọn theo xác thực của HAN và \val{mtv-m5-hit-mean} của metapath2vec.
Khi xếp hạng toàn bộ, nó đạt \val{gat-fhit-mean} so với \val{han-r6d2-fhit-s1} của HAN. Mức dẫn trước
của nó so với HAN lớn nhất với những người dùng có ít cạnh huấn luyện: \val{gat-lead-d1} với 1--5
cạnh, \val{gat-lead-d2} với 6--20 và \val{gat-lead-d3} từ 21 cạnh trở lên. Khác với hai
mô hình kia, GAT dựng mỗi người dùng từ các sản phẩm mà người đó đã nhận xét, trong khi HAN và
metapath2vec học một embedding tự do cho mỗi người dùng, embedding này chỉ nhận được ít tín hiệu từ
một vài cạnh. Điều này phù hợp với phát hiện của HGB rằng một GAT đồng nhất với đầu vào phù hợp là
một mô hình cơ sở mạnh \citep{lv2021hgb}. Điều này cũng gợi ý rằng ở đây cách biểu diễn người dùng
quan trọng hơn việc chọn những meta-path phía sản phẩm nào. Các meta-path phía người dùng vốn nằm
ngoài thiết kế của chúng tôi sẽ kiểm chứng được điều này.

HAN vượt metapath2vec về Hit@10 lấy mẫu trên cả
\val{han-wins-sampled} tập nhưng metapath2vec thắng khi xếp hạng toàn bộ ở \val{mtv-wins-full} tập,
nên thứ tự của hai mô hình phụ thuộc vào độ đo. Mọi mô hình đều vượt mức
sàn của mô hình phổ biến (\val{pop-hit} lấy mẫu, \val{pop-fhit} khi xếp hạng toàn bộ), trừ
metapath2vec trên M4 vốn đạt \val{mtv-m4-hit-s1}: tập này không có meta-path nào qua người
dùng nên người dùng không được huấn luyện.

\subsection{Độ nhạy}
\label{sec:res-sensitivity}

Với giai đoạn khởi động trọng số meta-path \val{wu5-warmup} epoch thay vì
\val{han-set-beta-warmup} (seed~1, cả 24 tập), HAN huấn luyện kém hơn về trung bình
(\val{wu5-diff-mean} Hit@10, từ \val{wu5-diff-min} đến \val{wu5-diff-max}), loại nhiều meta-path
hơn (\val{wu5-dropped} trọng số dưới \val{weight-drop} so với \val{han-dropped}) và sắp thứ tự các
tập khác đi (Spearman \val{wu5-rho} so với các lần chạy chính). Khi đó, các tập tìm được đứng
\val{wu5-4-pos} trong tám tập ở cả hai kích thước thay vì đứng cuối. Với $M = 6$ meta-path, phép so
sánh ở seed 1 đảo ngược: \code{a1} đạt \val{wu5-f6a1-hit} so với \val{wu5-m6-hit} của M6 vốn chỉ
giữ \val{wu5-m6-kept} trong $M = 6$ meta-path của nó. Phép kiểm tra này dùng một seed và không được lên
kế hoạch từ trước nên nó không thay đổi kết quả chính nhưng cho thấy phép so sánh phụ thuộc vào
cách HAN học trọng số meta-path: các kết luận của chúng tôi đúng với giai đoạn khởi động
\val{han-set-beta-warmup}~epoch.
